% HSC board-style CQ (সৃজনশীল) paper — compile: scripts/build.sh board-cq.tex
\documentclass[11pt]{article}
\usepackage[tiro]{edtech}
\geometry{left=20mm,right=24mm,top=14mm,bottom=16mm}
\pagestyle{empty}
\cqstyle{board}
\begin{document}
\boardcqheader{name=যমুনা, center=১৭৪, serial=১১২৬/২০২৬, set=০২ সেট,
  subject=পদার্থবিজ্ঞান, paper=প্রথম পত্র, time={২ ঘণ্টা ৩৫ মিনিট}, marks=৫০, code=১৭৪,
  note={যে কোনো পাঁচটি প্রশ্নের উত্তর দিতে হবে।}}

\cq{একজন ক্রিকেটার মাটি থেকে \SI{1.5}{m} উচ্চতায় একটি বলকে অনুভূমিকের সাথে \SI{30}{\degree} কোণে \SI{20}{m.s^{-1}} বেগে আঘাত করল। বলটি যে দিকে যাচ্ছে সেদিকে মাঠের সীমানা \SI{40}{m} দূরে এবং সেখানে একজন ফিল্ডার দাঁড়িয়ে আছে, যে লাফ দিয়ে সর্বোচ্চ \SI{2.5}{m} উচ্চতা পর্যন্ত বল ধরতে পারে। [$g=\SI{9.8}{m.s^{-2}}$]}
{প্রক্ষেপকের পাল্লা কাকে বলে?}
{সর্বাধিক উচ্চতায় প্রক্ষেপকের বেগ শূন্য হয় না কেন? ব্যাখ্যা কর।}
{বলটির বিচরণকাল নির্ণয় কর।}
{ফিল্ডার বলটি ধরতে পারবে কি না— গাণিতিক বিশ্লেষণের মাধ্যমে মতামত দাও।}

\cq{চিত্রে \SI{5}{kg} ভরের একটি বস্তুকে \SI{30}{\degree} আনত ঘর্ষণহীন তলের ওপর দিয়ে একটি দড়ির সাহায্যে টেনে তোলা হচ্ছে। দড়িটি আনত তলের সমান্তরাল এবং দড়ির টান \SI{40}{N}।\par
\begin{center}\begin{tikzpicture}[scale=0.9,font=\small]
 \fill[pattern=north east lines](0,0) rectangle (5.6,-0.18);
 \draw[thick](0,0)--(5.2,0)--(5.2,3)--cycle;
 \begin{scope}[shift={(2.6,1.5)},rotate=30]
   \draw[thick,fill=black!8](-0.45,0) rectangle (0.45,0.6); \node at (0,0.3){\SI{5}{kg}};
   \draw[-{Latex[length=2.5mm]},thick](0.45,0.3)--(1.8,0.3) node[above]{$T=\SI{40}{N}$};
 \end{scope}
 \draw (0.9,0) arc (0:30:0.9); \node at (1.4,0.28){\SI{30}{\degree}};
\end{tikzpicture}\end{center}}
{ঘাত বল কাকে বলে?}
{রকেটের গতি ভরবেগের সংরক্ষণ সূত্র মেনে চলে— ব্যাখ্যা কর।}
{বস্তুটির ত্বরণ নির্ণয় কর।}
{দড়ির টান অর্ধেক করা হলে বস্তুটি উপরে উঠবে নাকি নিচে নামবে? গাণিতিকভাবে বিশ্লেষণ কর।}

\cq{একটি গ্যাসের \SI{0}{\degreeCelsius} তাপমাত্রায় বিভিন্ন অবস্থায় চাপ ও আয়তন নিম্নরূপ:\par
\begin{center}\renewcommand\arraystretch{1.25}
\begin{tabular}{|c|c|c|c|c|}\hline
 চাপ (atm) & 0.35 & 0.50 & 0.65 & 0.85\\\hline
 আয়তন (L) & 3.80 & 2.66 & 2.05 & 1.56\\\hline
\end{tabular}\end{center}}
{আদর্শ গ্যাস কাকে বলে?}
{গ্যাসের অণুর গড় মুক্ত পথ চাপের ওপর কীভাবে নির্ভর করে? ব্যাখ্যা কর।}
{উদ্দীপকের গ্যাসটির মোল সংখ্যা নির্ণয় কর।}
{উদ্দীপকে উল্লিখিত গ্যাসটি গ্যাসের কোন সূত্রকে সমর্থন করবে? গাণিতিক যুক্তিসহ বিশ্লেষণ কর।}

\cq{পৃথিবীর ব্যাসার্ধ \SI{6400}{km} এবং ভূপৃষ্ঠে অভিকর্ষজ ত্বরণ \SI{9.8}{m.s^{-2}}। একটি কৃত্রিম উপগ্রহ ভূপৃষ্ঠ থেকে \SI{3600}{km} উচ্চতায় বৃত্তাকার কক্ষপথে পৃথিবীকে প্রদক্ষিণ করছে।}
{মুক্তিবেগ কাকে বলে?}
{পৃথিবীর কেন্দ্রে কোনো বস্তুর ওজন শূন্য হয় কেন? ব্যাখ্যা কর।}
{উপগ্রহটির কক্ষীয় বেগ নির্ণয় কর।}
{উপগ্রহটিকে ভূস্থির উপগ্রহে পরিণত করতে হলে তার উচ্চতার কীরূপ পরিবর্তন করতে হবে? গাণিতিকভাবে বিশ্লেষণ কর।}

\cq{একটি সরল ছন্দিত স্পন্দনে স্পন্দনশীল কণার সরণ $x=0.05\sin\left(4\pi t+\dfrac{\pi}{6}\right)$ m সমীকরণ দ্বারা প্রকাশিত। একটি সরল দোলকের দোলনকাল এই কণার পর্যায়কালের সমান।}
{পর্যায়কাল কাকে বলে?}
{সরল ছন্দিত স্পন্দনে কণার মোট শক্তি সংরক্ষিত থাকে কেন?}
{সরল দোলকটির কার্যকর দৈর্ঘ্য নির্ণয় কর।}
{$t=\SI{0.1}{s}$ সময়ে কণাটির গতিশক্তি ও বিভবশক্তির অনুপাত নির্ণয় করে মন্তব্য কর। [কণার ভর \SI{0.2}{kg}]}

\cq{চিত্রে দুটি ভিন্ন তারের জন্য পীড়ন বনাম বিকৃতি লেখচিত্র দেখানো হলো:\par
\begin{center}\begin{tikzpicture}[scale=0.85,font=\small]
 \draw[->](0,0)--(5,0) node[right]{বিকৃতি ($\times10^{-3}$)};
 \draw[->](0,0)--(0,3.6) node[above]{পীড়ন ($\times10^{8}$ N m$^{-2}$)};
 \draw[thick](0,0)--(3,3) node[right]{A};
 \draw[thick,dashed](0,0)--(4,2) node[right]{B};
 \foreach \x/\l in {1/1,2/2,3/3,4/4}{\draw(\x,0.06)--(\x,-0.06) node[below]{\l};}
 \foreach \y/\l in {1/1,2/2,3/3}{\draw(0.06,\y)--(-0.06,\y) node[left]{\l};}
\end{tikzpicture}\end{center}}
{হুকের সূত্রটি লেখ।}
{তাপমাত্রা বাড়লে তরলের সান্দ্রতা কমে কেন? ব্যাখ্যা কর।}
{A তারের ইয়ং-এর গুণাঙ্ক নির্ণয় কর।}
{একই দৈর্ঘ্য ও প্রস্থচ্ছেদের A ও B তারে সমান বল প্রয়োগ করলে কোনটিতে বেশি শক্তি সঞ্চিত হবে? গাণিতিকভাবে যাচাই কর।}

\cq{\SI{10}{kg} ভরের একটি বস্তু \SI{20}{m} উঁচু একটি দালানের ছাদ থেকে বিনা বাধায় নিচে পড়ছে। পড়ার সময় বস্তুটির গতিশক্তি ও বিভবশক্তি পর্যবেক্ষণ করা হলো।}
{ক্ষমতা কাকে বলে?}
{সংরক্ষণশীল বলের দ্বারা বদ্ধপথে কৃতকাজ শূন্য হয় কেন? ব্যাখ্যা কর।}
{ভূমি স্পর্শ করার মুহূর্তে বস্তুটির বেগ নির্ণয় কর।}
{পতনের কোন অবস্থানে বস্তুটির গতিশক্তি বিভবশক্তির তিনগুণ হবে? গাণিতিকভাবে বিশ্লেষণ কর।}

\cq{একটি চাকতির জড়তার ভ্রামক \SI{0.4}{kg.m^2}। চাকতিটি প্রতি মিনিটে \num{300} বার ঘুরছে। ঘূর্ণনরত অবস্থায় চাকতির উপর একটি ছোট বস্তু আস্তে রাখায় জড়তার ভ্রামক বেড়ে \SI{0.5}{kg.m^2} হলো।}
{কৌণিক ভরবেগ কাকে বলে?}
{ঘূর্ণন অক্ষের অবস্থান পরিবর্তনে জড়তার ভ্রামক পরিবর্তিত হয় কেন?}
{চাকতিটির প্রাথমিক ঘূর্ণন গতিশক্তি নির্ণয় কর।}
{বস্তুটি রাখার পর চাকতির ঘূর্ণন গতিশক্তি সংরক্ষিত থাকে কি না— গাণিতিকভাবে বিশ্লেষণ কর।}
\end{document}
